\section{The adversary}\label{sec:algorithm}

The adversary maintains a rooted tree of groups. A group records comparisons
among items whose relative order is still being resolved there. An item may
leave a group by fixing its position and continuing in a child at that
position. We call this operation a \emph{refinement}.

Fix an integer $k\ge16$. The algorithm uses the parameters
\begin{equation}\label{eq:parameters}
\eta=\frac1{100k},\qquad H=2^{2k},\qquad
B=(k+1)^{2H+1},\qquad R=2^{2kB}.
\end{equation}
Between comparison calls, the graph in each group has degree less than $k$
and no directed path on $H$ vertices. These two restrictions will make the
probability calculations local. The threshold $1/R$ determines when a queried
coordinate must first be refined. Section~\ref{sec:runtime} selects $k$ as a
function of $n$; for now it is fixed.

\subsection{The state of a group}

At each node $v$, every item represented there is either \emph{free} or belongs
to a child subtree. A free item $x$ has an auxiliary coordinate $X_x$ in an open
interval $(a_x,b_x)\subseteq(0,1)$. These coordinates describe possible orders;
they are not the unknown key values. A directed acyclic graph on the free
items stores inequalities $X_x<X_y$. Each child has a fixed position
$c\in(0,1)$, its \emph{anchor}, and its own copy of the unit interval.
Different children of one node have different anchors.

Each original item has exactly one free coordinate in the entire tree. At
every ancestor of its node it is represented by the anchor of the child
containing it. Initially all $n$ items are free at the root, their intervals
are $(0,1)$, and there are no arrows or children.

The feasible region of a node is
\begin{equation}\label{eq:node-region}
K_v=\{X:a_x<X_x<b_x\text{ for all free }x,
             \quad X_x<X_y\text{ for every arrow }x\to y\}.
\end{equation}
Write $V_v=\vol(K_v)$, with volume one when the node has no free items, and
define
\begin{equation}\label{eq:Z}
Z=\prod_v V_v.
\end{equation}
Thus $Z=1$ initially. We shall prove that every reached state satisfies
$0<Z\le1$. The algorithm never
calculates this product. It calculates probabilities only in small induced
subgraphs, as specified in Section~\ref{sec:oracle}.

Only actual comparisons create arrows. The graph is simple: repeated
comparisons do not create parallel arrows, and transitive consequences are
not inserted. The number of stored arrows will matter in the proof.

We maintain all interval bounds implied by the arrows. Write $y\preceq x$
if $y=x$ or a directed path leads from $y$ to $x$. The \emph{effective interval}
of $x$ is $(\ell_x,r_x)$, where
\begin{equation}\label{eq:effective}
\ell_x=\max_{y\preceq x}a_y,\qquad
r_x=\min_{y\succeq x}b_y,\qquad W_x=r_x-\ell_x.
\end{equation}
The operation \textnormal{\textsc{Tighten}} replaces every stored interval by its effective
interval. It is performed after every change to an interval or arrow, before
testing widths or calculating probabilities. The resulting bounds remain
stored when arrows are later removed. A coordinate is \emph{wide} if
$W_x\ge1/R$, and \emph{narrow} otherwise.

\subsection{Choosing an answer}

To compare original items $A$ and $B$, descend while both belong to the same
child. At the first node where their representations differ, each is either
a free coordinate or a child anchor. Two different anchors determine the
answer. Otherwise at least one representation is free.

Suppose first that both are free. Under the uniform volume measure on $K_v$,
let
\[
p=\Prb(X_A<X_B)=\frac{\vol(K_v\cap\{X_A<X_B\})}{V_v}.
\]
The probability notation is only a ratio of volumes; the adversary makes no
random choices. Given an estimate $\widehat p$ with
$|\widehat p-p|\le\eta$, it answers $A<B$ if $\widehat p\ge1/2$, and $B<A$
otherwise. Since equality has volume zero, the chosen answer retains at
least $1/2-\eta$ of $V_v$. It is stored as an arrow. For a free coordinate
and an anchor $c$, write $x$ for the free item and use $\Prb(X_x<c)$.
Store the resulting upper or lower interval bound, reversing the reported
label order when the free item is $B$.

The local estimator requires a wide coordinate for a threshold query, and
at least one wide endpoint for a two-coordinate query. If this condition
fails, the adversary refines a queried item before answering.

\subsection{Choosing an anchor and refining an item}

Every anchor is the midpoint of a dyadic cell:
\begin{equation}\label{eq:dyadic}
c=(j+\tfrac12)2^{-m},\qquad 0\le j<2^m,\quad m\ge0.
\end{equation}
We call $m$ its \emph{address length}: $m$ binary digits specify the cell.
The reduced denominator of $c$ is $2^{m+1}$, so its position determines
$m$ uniquely. Reusing an anchor therefore also reuses its address length.

There are two rules for choosing $c$.

\begin{enumerate}
\item For a wide item, start with the cell $[0,1]$. At its midpoint $c$,
estimate $F_x(c)=\Prb(X_x<c)$ to error at most $\eta$.
If the estimate is below $1/4$, retain the right half; if it is above $3/4$,
retain the left half; otherwise accept $c$. Repeat within the retained cell.
\item For a narrow item, choose the dyadic length $b=2^{-m}$ satisfying
$W_x/8<b\le W_x/4$. Choose the midpoint of a cell of this length nearest
to $(\ell_x+r_x)/2$, taking the smaller midpoint in case of a tie.
\end{enumerate}

Both choices lie strictly inside the effective interval. Their density
guarantees are proved in Section~\ref{sec:oracle}.

To refine $x$ at $c$, fix $X_x=c$, remove $x$ from the parent's free graph,
and replace each incident arrow by the corresponding bound on its other
endpoint. Thus $y\to x$ becomes $X_y<c$, and $x\to z$ becomes $X_z>c$.
Retain all other stored bounds and tighten the parent. Then admit $x$ to
the child at $c$, creating the child if necessary, as an isolated coordinate
with interval $(0,1)$. The child's previous data stay unchanged.

\begin{figure}[tb]
\centering
\begin{tikzpicture}[>=Latex,font=\small,
  item/.style={circle,draw=ink,fill=white,minimum size=6mm,inner sep=1pt},
  group/.style={draw=ink!65,rounded corners=2pt,fill=pale}]
\begin{scope}[shift={(-3.4,0)}]
\draw[group] (-2.7,-.9) rectangle (2.7,1.15);
\node at (0,.8) {Before refinement};
\node[item] (y) at (-1.7,0) {$y$};
\node[item] (x) at (0,0) {$x$};
\node[item] (z) at (1.7,0) {$z$};
\draw[->,thick] (y)--(x);
\draw[->,thick] (x)--(z);
\end{scope}
\draw[->,thick] (-.45,.1)--(.45,.1);
\begin{scope}[shift={(3.4,0)}]
\draw[group] (-2.7,-.9) rectangle (2.7,1.15);
\node at (0,.8) {After fixing $X_x=c$};
\draw[thick] (-2.35,-.15)--(2.35,-.15);
\fill (0,-.15) circle (2pt);
\node[below] at (0,-.18) {$c$};
\node at (-1.5,.15) {$X_y<c$};
\node at (1.5,.15) {$c<X_z$};
\draw[group] (-2.35,-2.7) rectangle (2.35,-1.6);
\node at (0,-1.9) {Child at $c$};
\node at (0,-2.35) {old items, plus $X_x\in(0,1)$};
\draw[->,thick,accent] (0,-.9)--(0,-1.6);
\end{scope}
\end{tikzpicture}
\caption{A refinement preserves the comparisons incident to $x$ as interval
bounds in the parent. The relative order of $x$ and the items already at the
same anchor is resolved inside the child.}
\label{fig:refinement}
\end{figure}

\subsection{The complete comparison procedure}

The following procedures specify all state changes. Ties between eligible
items are resolved by label order. A path is selected by the lexicographically
first eligible label sequence. Degree counts neighbors without regard to
edge direction; path length below counts vertices.

\begin{quote}
\textbf{Restore$(v)$.}
\begin{enumerate}
\item While the free graph of $v$ has a vertex of degree at least $k$,
refine such a vertex.
\item While it contains a directed path on $H$ vertices, select one such
path and refine all its vertices in label order. Complete this entire list
before selecting another path.
\end{enumerate}
\end{quote}

\begin{quote}
\begin{minipage}{\linewidth}
\textbf{Compare$(A,B)$.}
\begin{enumerate}
\item Descend through common children to the first node $v$ where the
representations differ. If both are anchors, return their order.
\item Tighten $v$. If none of the free queried coordinates is wide,
refine one of them and restart step 1.
\item Estimate the probability of $A<B$ to additive error $\eta$ and select
an answer by the rule above. Store its arrow or interval bound.
\item Tighten $v$, run \textnormal{\textsc{Restore}}$(v)$, and return the selected answer.
\end{enumerate}
\end{minipage}
\end{quote}

Every answer is also entered in a separate comparison graph on the original
labels, before restoration or, for an anchor answer, before returning.
This graph records the public transcript. Its edges do not count toward the
degree restrictions or the active arrow count used below. It will supply a
compatible total order for the finite stopping rule in
Section~\ref{sec:runtime}.
